\documentclass[dvipsnames]{article}

\usepackage[utf8]{inputenc}
\usepackage{fontspec}
\usepackage{polyglossia}
\usepackage{unicode-math}

\setmainfont{CMU Serif}
\setsansfont{CMU Sans Serif}
\setmonofont{CMU Typewriter Text}
% \setmathfont{Latin Modern Math}

\usepackage[letterpaper,top=2cm,bottom=2cm,left=3cm,right=3cm,marginparwidth=1.75cm]{geometry}

\usepackage{amsmath}
\usepackage{fancyhdr}
\usepackage[colorlinks=true, allcolors=blue]{hyperref}
\usepackage[backend=biber]{biblatex}

\usepackage{minted2}
\usemintedstyle{xcode}
\setminted{escapeinside=||}
\newmintinline[python]{python}{fontsize=\small}
\newmintinline[code]{c}{fontsize=\small}
\let\codeOrig\code
\renewcommand{\code}[2][]{%
  \ifmmode
    \if\relax\detokenize{#1}\relax
      \mathchoice{\mbox{\normalfont\ttfamily #2}}{\mbox{\normalfont\ttfamily #2}}{\mbox{\normalfont\ttfamily #2}}{\mbox{\normalfont\ttfamily #2}}%
    \else
      \mathchoice{\mbox{\normalfont\footnotesize\ttfamily #2}}{\mbox{\normalfont\footnotesize\ttfamily #2}}{\mbox{\normalfont\footnotesize\ttfamily #2}}{\mbox{\normalfont\footnotesize\ttfamily #2}}%
    \fi
  \else
    \codeOrig[#1]{#2}%
  \fi}

\usepackage{mdframed}

\title{Языки программирования. Семантика и системы типов}
\author{Николай Кудасов}
\date{16 февраля 2026}

\addbibresource{../references.bib}

\begin{document}

\pagestyle{fancy}
\fancyhf{}
\chead{Задание выполнил \textcolor{blue}{Павел Григоренко}}
\cfoot{Языки программирования. Семантика и системы типов (весна 2026, Университет ИТМО)}

\section*{Тема 1. Теоретическое задание \textcolor{blue}{(выполнил Павел Григоренко)}}%

\begin{enumerate}
  \item
В бестиповом $\lambda$-исчислении~\cite[\S 5]{TaPL-ru} выполните задания по подстановке с избежанием захвата, строгому вычислению и безымянному представлению.

  \begin{enumerate}
  \item
  Выполните подстановку \code{[y |$\mapsto$| (λa. x a)(λb. z b)](λx. λw. λz. y x z)}. \\
  При необходимости выполните $\alpha$-конверсию (переименование связанных переменных) и укажите итоговый результат подстановки.

  \item
  Переведите следующий терм в безымянное представление (индексы де Брёйна)~\cite[\S 6]{TaPL-ru} и
  вычислите его в безымянной форме, используя \emph{строгое} вычисление (один шаг — один переход ($\longrightarrow$) по операционной семантике).
  \begin{mdframed}
    \code{(λf. λx. f (f (f x))) (λg. λy. g (g y)) (λz. z)}
  \end{mdframed}
  \end{enumerate}

  \color{blue}
  \textbf{Решение.}
  \color{black}

  \begin{enumerate}
  \item

  \code{[y |$\mapsto$| (λa. x a)(λb. z b)](λx. λw. λz. y x z)}

  \code{[y |$\mapsto$| (λa. x a)(λb. z b)](λx. λw. λt. y x t)}

  \code{[y |$\mapsto$| (λa. x a)(λb. z b)](λu. λw. λt. y u t)}

  \code{λu. λw. λt. (λa. x a)(λb. z b) u t}

  \item

  \code{(λf. λx. f (f (f x))) (λg. λy. g (g y)) (λz. z)}

  \code{(λ λ 1 (1 (1 0))) (λ λ 1 (1 0)) (λ 0)}

  \code{(λ (λ λ 1 (1 0)) ((λ λ 1 (1 0)) ((λ λ 1 (1 0)) 0))) (λ 0)}

  \code{(λ λ 1 (1 0)) ((λ λ 1 (1 0)) ((λ λ 1 (1 0)) (λ 0)))}

  \code{(λ λ 1 (1 0)) ((λ λ 1 (1 0)) (λ (λ 0) ((λ 0) 0)))}

  \code{(λ λ 1 (1 0)) ((λ λ 1 (1 0)) (λ (λ 0) 0))}

  \code{(λ λ 1 (1 0)) ((λ λ 1 (1 0)) (λ 0))}

  \code{(λ λ 1 (1 0)) (λ (λ 0) ((λ 0) 0))}

  \code{(λ λ 1 (1 0)) (λ (λ 0) 0)}

  \code{(λ λ 1 (1 0)) (λ 0)}

  \code{λ (λ 0) ((λ 0) 0)}

  \code{λ (λ 0) 0}

  \code{λ 0}

  \end{enumerate}


  \clearpage
  \item Для бестипового $\lambda$-исчисления с булевыми и арифметическими выражениями~\cite[\S 5]{TaPL-ru} докажите по \emph{индукции по правилам вывода} одношагового вычисления, что при вычислении число свободных переменных не увеличивается:
  \[
    \text{если } t \longrightarrow t', \text{ то } FV(t') \subseteq FV(t) \quad \text{(в частности, } |FV(t')| \leq |FV(t)|\text{).}
  \]

  \color{blue}
  \textbf{Решение.}
  \color{black}

  \begin{itemize}
    \item Набор свободных переменных в if это объединение наборов свободных переменных
    его условия, и двух веток.
    \begin{itemize}
      \item При выводе в одну из веток при явном значении условия, остаёмся лишь с набором свободных переменных этой ветки. Очевидно, что подмножество.
      \item При выводе, когда условие выводится в новую форму,
      (по индуктивной гипотезе) множество его свободных переменных становится нестрогим подмножеством того, что было раньше. Очевидно, что если его теперь объединит
      с множествами для веток, получим подмножество того, что было раннее.
    \end{itemize}
    \item $\code{pred 0}\mapsto\code{0}$: было и осталось $\emptyset$.
    \item $\code{pred (succ v)}\mapsto\code{v}$: было и осталось $FV(v)$.
    \item $\code{iszero 0}\mapsto\code{true}$: было и осталось $\emptyset$.
    \item $\code{iszero (succ v)}\mapsto\code{false}$: $\emptyset\subseteq FV(v)$.
    \item Правила вывода \code{pred}, \code{succ} и \code{iszero} при выводе их аргумента
    очевидно подходят по индуктивной гипотезе.
    \item Правила вывода аппликации при выводе её левой или её правой части тоже очевидно подходят по индуктивной гипотезе.
    \item Вычисление лямбды подстановкой аргумента. Будем считать, что ни одна абстракция в левой части не использует имена каких-либо свободных переменных правой части. Иначе выполним $\alpha$ преобразование. Теперь все новые вхождения правой части в итоговое выражение
    очевидно не поменяли своё множество свободных переменных. Что касается левой части,
    в её теле мы убрали все вхождения той переменной, по которой шла абстракция и саму абстракцию по этой переменной. Так что по определению свободных переменных в абстракции тоже очевидно нет изменений.
  \end{itemize}

  \clearpage
  \item Приведите терм $\code{t}$ в бестиповом $\lambda$-исчислении с арифметическими выражениями~\cite[\S 5]{TaPL-ru}, такой что
  \begin{itemize}
    \item размер терма $\mathsf{size}(\code{t}) \geq 20$
    \item $\code{t (λx. λy. x)} \longrightarrow^* \code{succ 0}$
    \item $\code{t (λx. λy. y)} \longrightarrow^* \code{succ (succ 0)}$
  \end{itemize}

  \color{blue}
  \textbf{Решение.}
  \color{black}

  \code{t := λf. (pred (succ (pred (succ (pred (succ (pred (succ (pred (succ (pred (succ}
  \\
  \code{(pred (succ (pred (succ (pred (succ (pred (succ (f (succ 0) (succ (succ 0)))))))))))))))))))))))}

  \clearpage
  \item Преобразуйте фрагмент программы на Python ниже в эквивалентный фрагмент, применив операцию инлайнинга функций \python{make_reporter} и \python{report_cost}.
  Если для корректного преобразования требуется переименование переменных или другие изменения — явно дайте пояснение этим действиям.
  Убедитесь в корректности преобразования проверкой в интерпретаторе.

\begin{minted}[linenos,frame=single,fontsize=\footnotesize]{python}
def make_reporter(level):
  threshold = level * 2
  def report_cost(msg_id):
    if msg_id >= 2:
      return (level * msg_id) + threshold
    return level * msg_id
  return report_cost

def total_diagnostic_cost():
  levels = list(map(int, input().split()))
  report_order = list(map(int, input().split()))
  msg_id = int(input())

  reporters = {}
  for level in levels:
    reporters[level] = make_reporter(level)

  total = 0
  for lev in report_order:
    total += reporters[lev](msg_id)

  return total

total_diagnostic_cost()
\end{minted}

  \color{blue}
  \textbf{Решение.}
  \color{black}

  Переименовываем аргумент у \code{report_cost} из \code{msg_id} в \code{m}.
  Меняем её две ветки на тернарник, и превращаем саму функцию в лямбду.
  Подставляем \code{report_cost} на место её вызова.
  Меняем возврат значения на запись его туда, куда раньше оно записывалось после возврата.

\begin{minted}[linenos,frame=single,fontsize=\footnotesize]{python}
def total_diagnostic_cost():
  levels = list(map(int, input().split()))
  report_order = list(map(int, input().split()))
  msg_id = int(input())

  reporters = {}
  for level in levels:
    threshold = level * 2
    reporters[level] = lambda m: (level * m) + threshold if m >= 2 else level * m

  total = 0
  for lev in report_order:
    total += reporters[lev](msg_id)

  return total

total_diagnostic_cost()
\end{minted}

\end{enumerate}

\clearpage

\section*{Список литературы}
\printbibliography[heading=none]

\end{document}
